\documentclass[11pt]{article}

\usepackage[a4paper,margin=0.84in]{geometry}
\usepackage[T1]{fontenc}
\usepackage[utf8]{inputenc}
\usepackage{lmodern}
\usepackage{amsmath,amssymb,amsthm,mathtools}
\usepackage{booktabs,array,tabularx,microtype,xcolor}
\usepackage[numbers,sort&compress]{natbib}
\usepackage[colorlinks=true,linkcolor=blue!65!black,
  citecolor=blue!65!black,urlcolor=blue!70!black]{hyperref}
\usepackage[nameinlink,noabbrev]{cleveref}
\setlength{\emergencystretch}{3em}

\hypersetup{
  pdftitle={Actual Active Support on an Empty Eligible Family},
  pdfauthor={Liang Wang},
  pdfsubject={A vacuity firewall for the H1 carrier root},
  pdfkeywords={active support, physical coefficient, empty domain, vacuity, H1}
}

\newtheorem{theorem}{Theorem}[section]
\newtheorem{proposition}[theorem]{Proposition}
\theoremstyle{definition}
\newtheorem{definition}[theorem]{Definition}
\theoremstyle{remark}
\newtheorem{remark}[theorem]{Remark}
\newcommand{\Lzero}{\mathrm{L0}}
\newcommand{\Lone}{\mathrm{L1}}
\newcommand{\Ltwo}{\mathrm{L2}}
\newcommand{\NT}{\textnormal{\textsc{not-testable}}}
\newcolumntype{Y}{>{\raggedright\arraybackslash}X}

\title{\textbf{Actual Active Support on an Empty Eligible Family:}\\
\textbf{A Vacuity Firewall for the H1 Carrier Root}}
\author{Liang Wang\\
School of Artificial Intelligence and Automation,\\
Huazhong University of Science and Technology, Wuhan 430074, P.R. China\\
\texttt{wangliang.f@gmail.com}}
\date{July 2026}

\begin{document}
\maketitle

\begin{abstract}
The H1 occurrence architecture requires an actual active-support
certificate independently of a source-backed local occurrence-edge
family.  TPC-175 and TPC-176 leave no eligible production carrier in
the frozen TPC-133--172 source scope.  We audit literal physical
coefficients only on that eligible domain.  No carrier is tested, so
the counts of nonzero, zero, and missing coefficients are all zero.

The resulting universal support statement is vacuous and cannot
supply an existential witness.  We formalize this as a firewall:
the empty-domain audit passes, while
\(\mathrm{H1.actual\_active\_support\_certificate}\) remains
\(\NT\).  The independent H9 literal-weight registry is neither
imported nor closed and retains decay axis \texttt{NONE}.  This is an
\(\Lone\) interface obstruction, not physical arithmetic progress or
program-positive \(\Ltwo\).
\end{abstract}

\section{Two independent prerequisites}

TPC-166 splits the former occurrence crosswalk monolith into three
independent roots \cite{WangTPC166}.  Write
\[
\begin{aligned}
 E&=\text{a source-backed production local edge family},\\
 A&=\text{actual active support of its physical carriers},\\
 M&=\text{a canonical or minimal physical representation}.
\end{aligned}                                                     \tag{1}
\]
Neither \(E\Rightarrow A\) nor \(A\Rightarrow E\) is available.
TPC-177 studies \(A\) without collapsing this independence.

For a candidate carrier \(c\), an active-support record must contain
at least a literal target-expression locator, the literal physical
coefficient \(w(c)\), its target normalization, and a proof that
\[
 w(c)\ne0.                                                        \tag{2}
\]
Formal appearance in a schema, archive, completion fiber, or
synthetic fixture does not satisfy (2).

\section{The eligible domain}

TPC-175 certifies no qualifying production local edge in the frozen
source inventory, and TPC-176 therefore exposes no eligible carrier:
\[
 \mathcal C_{\rm elig}=\varnothing.                               \tag{3}
\]
The support audit is deliberately restricted to (3).  Adding a
formal archived cut or a synthetic carrier would fabricate the
missing source-backed edge.

Define
\[
\begin{aligned}
 \mathcal C_{+}&=\{c\in\mathcal C_{\rm elig}:w(c)\ne0
                     \text{ is source-proved}\},\\
 \mathcal C_{0}&=\{c\in\mathcal C_{\rm elig}:w(c)=0
                     \text{ is source-proved}\},\\
 \mathcal C_{?}&=\{c\in\mathcal C_{\rm elig}:
                     \text{no complete coefficient record exists}\}.
\end{aligned}                                                     \tag{4}
\]
These three classes partition the eligible domain when every carrier
is assigned exactly one audit outcome.

\begin{proposition}[Exact coefficient ledger]
For the TPC-176 domain,
\[
 (|\mathcal C_{\rm elig}|,|\mathcal C_{+}|,
   |\mathcal C_{0}|,|\mathcal C_{?}|)=(0,0,0,0).                  \tag{5}
\]
\end{proposition}

\begin{proof}
Immediate from (3) and the definitions in (4).
\end{proof}

\section{The vacuity firewall}

On the empty domain, the sentence
\[
  \forall c\in\mathcal C_{\rm elig},\quad w(c)\ne0                \tag{6}
\]
is true in classical logic.  But the H1 root requires production
content:
\[
  \exists c\in\mathcal C_{\rm elig}
  \quad\text{with a source-backed local edge and }w(c)\ne0.       \tag{7}
\]

\begin{theorem}[TPC-177 vacuity firewall]
Equation (6) does not prove (7).  The executable audit status is
\texttt{VACUOUS\_EMPTY\_ELIGIBLE\_DOMAIN}, and the H1
actual-active-support root remains \(\NT\).
\end{theorem}

\begin{proof}
An existential witness for (7) would be an element of the empty set
in (3), which is impossible within the declared audit instance.
This does not prove that an actual active carrier is mathematically
impossible; it proves that the current source-locked input supplies
none.  Hence neither \texttt{PROVED} nor a global
\texttt{STOP\_ROUTE} classification is valid.
\end{proof}

The distinction may be summarized as follows.

\begin{center}
\begin{tabularx}{0.96\linewidth}{@{}l c Y@{}}
\toprule
Statement & Value & Evidential consequence\\
\midrule
empty-domain audit internally consistent & true &
the ledger and partition checks pass\\
all eligible carriers active & vacuously true &
no existential content\\
an active production carrier exists & unproved &
H1 support root remains \(\NT\)\\
active carriers do not exist mathematically & unproved &
frozen-corpus scope cannot imply nonexistence\\
\bottomrule
\end{tabularx}
\end{center}

\section{Separation from the H9 literal-weight registry}

The H1 support root and
\(\mathrm{H9.literal\_weight\_registry}\) have different jobs.  H1
asks whether a source-backed occurrence carrier is present with a
nonzero physical coefficient.  H9 supplies literal weight data to a
later arithmetic synthesis.  A registry has data semantics and
decay axis \texttt{NONE}; it does not manufacture a cancellation
exponent.

TPC-177 neither scans nor imports the H9 registry.  Thus it records
\[
\begin{aligned}
 \text{H1 active-support status}&=\mathsf{NOT\_TESTABLE},\\
 \text{H9 literal-weight status}&=\mathsf{NOT\_IMPORTED},\\
 \text{decay created by either registry}&=0.
\end{aligned}                                                     \tag{8}
\]
Even if a future H9 registry names a literal coefficient, an
independent source-backed H1 carrier would still be required.

\section{Reproducibility and claim boundary}

The standard-library audit source-locks the TPC-175 maximal family
and TPC-176 coverage ledger, recomputes (3)--(5), and checks the
existential firewall.  Mutations that fabricate an eligible carrier,
promote vacuous truth to an H1 certificate, import H9 implicitly,
assign H9 a decay exponent, or promote the result to \(\Ltwo\) are
rejected.  Hashes have integrity semantics only.

No fixed physical \(h_0=2\), named phase, deterministic endpoint,
target normalization, or endpoint slack is consumed.  Accordingly
there is no strict \(1/400\) gain, prime-pair lower bound, or
twin-prime conclusion.

\section{Conclusion}

An empty eligible domain makes an audit easy but a theorem
impossible: the universal support sentence is vacuous, and the
required existential production witness is absent.  TPC-177 keeps
the H1 support root open and keeps it disjoint from the H9
literal-weight registry.  TPC-178 applies the same eligibility
discipline to canonical and minimal physical representations.

\nocite{*}
\bibliographystyle{plainnat}
\bibliography{references}

\end{document}
